\tocsection{\textbf{9/17/26 Lecture 7: }Exponential Families}

Recall, an exponential family looks like $\cal{P} = \{p_{\eta} \col \eta \in \Xi\}$ with
\[ p_{\eta}(x) = e^{\eta'T(x) - A(\eta)}h(x) \]
These give nice forms for moments, which are the differential identities
\[ \nabla A(\eta) = \E_{\eta}[T(X)], \quad \nabla^2 A(\eta) = \Var_{\eta}(T(X)), M_{\eta}^{T(X)}(u) = e^{A(\eta + u) - A(\eta)} \]
The reason these are nice is that a product of exponential family is an exponential family, so these moments are easy to compute.\\

Also, one last note, we may also find it convenient to parameterize our exponential families in a different form, like
\[ p_{\eta}(x) = e^{\eta(\theta)'T(x) -B(\theta)}h(x) \]
where $B(\theta) = A(\eta(\theta))$, so we may have to ``massage'' these forms into the canonical form. Here are some examples.

\begin{example}{\textbf{Example}}
    Let $X \sim \cal{N}(\mu, \sigma^2)$ with $\mu \in \R$ and $\sigma^2 > 0$ unknown. Then, we see that when $\theta = (\mu, \sigma)$, we have
    \[ p_{\theta}(x) = \frac{1}{\sqrt{2\pi \sigma^2}}e^{-\frac{1}{2\sigma^2}(x - \mu)^2} \]
    but by expanding, we can write this like
    \[ p_{\theta}(x) = e^{\frac{\mu}{\sigma^2}x - \frac{x^2}{2\sigma^2} - \frac{\mu^2}{2\sigma^2} - \frac{1}{2}\log(2\pi \sigma^2)} \]
    In this case, we see that
    \[ \eta(\theta) = \begin{bmatrix} \frac{\mu}{\sigma^2} \\ -\frac{1}{2\sigma^2}\end{bmatrix}, \quad T(x) = \begin{bmatrix} x \\ x^2 \end{bmatrix}, \quad B(\theta) = \frac{\mu^2}{2\sigma^2} + \frac{1}{2}\log(2\pi\sigma^2), \quad h(x) = 1 \]
    where we can put it back in canonical form writing
    \[ A(\eta(x)) = \frac{-\eta_1^2}{4\eta_2} + \frac{1}{2}\log(-\frac{\pi}{\eta_2}) \]
    so
    \[ p_{\eta}(x) = e^{\eta'\begin{pmatrix} x \\ x^2\end{pmatrix} - A(\eta)}h(x) \]

    Notice something interesting. $T(X)$, the sufficient statistic, was supposed to be a data reduction, but now we turned it into a vector that has arguably (not really) more data! Where this helps however is when you have multiple variables; then, the sufficient statistic is still just a vector whereas the data is a $n$-tuple. 
\end{example}

\begin{example}{\textbf{Example}}
    Let $X \sim \Binom(n, \theta)$ with $\theta$ unknown, then, we can write
    \begin{align*}
        p_{\theta}(x) &= \theta^x(1 - \theta)^{n - x}{n \choose x}
    \end{align*}
    Now, it seems like we will need a sufficient statistic consisting of $x$ and $n - x$, but notice that these have a very intimate relation; their sum must be $n$! Thus, we can reduce the dimension by writing
    \begin{align*}
        p_{\theta}(x) &= \left(\frac{\theta}{1 - \theta}\right)^{x}(1 - \theta)^{n}{n \choose x}\\
        &= e^{x\log\frac{\theta}{1-\theta} - n\log\frac{1}{1 - \theta}}{n \choose x}
    \end{align*}
    so we have
    \[ \eta = \log\frac{\theta}{1 - \theta},\quad x = T(x),\quad A(\eta) = n\log\frac{1}{1 - \theta},\quad h(x) = {n\choose x} \]
\end{example}

One way we can think of exponential families is as ``exponential tilting'' of the carrier $h(x)$. That is, we start with our carrier distribution $h(x)$, then we tilt by $e^{\eta'T(x)}$, and then we renormalize using $e^{-A(\eta)}$. What this does is that it tilts the density towards points for which $\eta'T(x)$ is large, and away from points where $\eta'T(x)$ is small.\\

For example, consider the exponential family
\[ p_{\eta}(x) \propto e^{\eta_1x + \eta_2x^2 + \eta_3\cos(100x)}(1 - \abs{x}) \]
Increasing $\eta_1$ pushes the distribution to the right, increasing $\eta_2$ pushing the distribution to the edges, and increasing $\eta_3$ makes $p_{\eta}$ oscillate ``like cosine''.\\

Also, note that decomposition of an exponential family into $\eta, T, A, h$ is very not unique. Firstly, it only matters what the span of $T(x)$ is, as choosing $\eta$ allows us to map any element of the span of $T$ to another. Additionally, we can always add constants to the terms in $T(x)$, because that can be taken out and absorbed into the $A(\eta)$ term.\\

Now, we will prove a quick result which will reinforce the notion that the sufficient statistic is what we should care about.

\begin{lemma}[\textbf{Proposition 7.1}]
    Suppose $X \sim p_{\eta}(x) = e^{\eta'T(x) - A(\eta)}$ with respect to $\mu$. Then, 
    \[ T(x) \sim q_{\eta}(t) = e^{\eta't - A(\eta)} \]
    with respect to $\nu = \mu \circ T^{-1}$. 
\end{lemma}
\begin{proof}{Proof:}
    We see that
    \begin{align*}
        \bb{P}_{\eta}(T(X) \in B) &= \int_{\chi} 1_{B}(T(x))e^{\eta'T(x) - A(\eta)}d\mu(x)\\
        &= \int_{t} 1_{B}(t)e^{\mu't - A(\eta)}d\nu(t)
    \end{align*}
    where we are basically doing a change of coordinates, $x \to T(x)$. This can be easier seen with the counting measure
    \[ \sum_{x \col T(x) \in B} e^{\eta'T(x) - A(\eta)}\mu(\{x\}) = \sum_{t \in B} e^{\eta't - A(\eta)}\mu(T^{-1}(t)) \]
    In words, $p_{\eta}(x)$ is a density for each $x$, $q_{\eta}(t)$ is a density for each slice $\{x \col T(x) = t\}$. 
\end{proof}

Through these lens, exponential tilt is taking a carrier density and tilting it on the slices $\{x \col T(x) = t\}$. This makes sense to do because tilting by $e^{\eta'T(x)}$ does not depend on $x$, but on where $T$ sends $x$. Also, this is nice because if $X$ is an exponential family, $T(X)$ is also an exponential family, but it actually is of a special kind.

\begin{definition}[\textbf{Definition 7.2:} Canonical Form]
    If the densities of our exponential family are of the form
    \[ p_{\eta}(x) = e^{\eta'x - A(\eta)} \]
    we say this exponential family is in \emph{canonical form.} 
\end{definition}

From the remarks above, notice we can always put an exponential family in canonical form by absorbing $h$ into our measure, and performing a sufficiency reduction $x \to T(x)$.\\

\subsection*{Minimal Form}
\sectionspace

\begin{definition}[\textbf{Definition 7.3:} Minimal Form]
    We say an exponential family
    \[ p_{\eta}(x) = e^{\eta'T(x) - A(\eta)}h(x) \]
    is in \emph{minimal form} if $\eta \in \Xi$ and $T(x)$ satisfy no linear constraints. That is, there are no $a \ne 0$ and $b \in \R$ such that $\eta'a = b$ or $T(x)'a = b$ $\cal{P}$ almost surely. (where $\cal{P}$ almost surely means this holds on a set of full measure for some $\eta$). 
\end{definition}
The reason we care about minimal form is because, if an exponential family is not in minimal form, we can use the linear constrains to reduce the dimension from a $s$-dimensional exponential form to an $r$-dimensional exponential form with $r < s$. 

\begin{lemma}[\textbf{Lemma 7.4}]
    If $p_{\eta}$ is minimal (in minimal form), then $T(x)$ is minimal sufficient.
\end{lemma}
\begin{proof}{Proof:}
    We want to show that if $\ell(\cdot, x) = \ell(\cdot, y)$ up to scalar, then $T(x) = T(y)$. Notice that since $T$ is sufficient, the converse of this statement holds.\\

    First, we see that
    \[ \ell(\eta, x) - \ell(\eta, y) = \eta'(T(x) - T(y)) + \log(h(x)) - \log(h(y)) \]
    letting $a = T(x) - T(y)$, then if $\ell(\eta, x) - \ell(\eta, y) = C$ for a constant $C$, then we see that
    \[ C - \log(h(x) - \log(h(y))) = \eta'a \]
    Then, as the LHS is constant in $\eta$, then if $a \ne 0$, this contradicts the fact that $p_{\eta}$ is minimal. Thus, we must have that $a = 0$, so $T(x) = T(y)$. 

    $\qed$
\end{proof}

Consider the following diagram

\fig{0.8}{Lecture7Fig1}{}
where $(A), (B)$ and $(C)$ are exponential families living inside $\Xi_{1}$. Which ones are minimal? Firstly, we can definitely see that $A$ and $B$ are not minimal, because there is no way to write the $\eta$ that belong to those families as a solution to a linear equation. However, for $(C)$, we can.\\

that is, what we can do is pick a base point $\eta_0$ on that line, and for any other $\eta$ on that line, we can write $\eta = \eta_0 + \theta \gamma$ for $\gamma$ the unit vector in the direction of the line. Then, we see that we have
\begin{align*}
    e^{\eta'T(x) - A(\eta)}h(x) &= e^{(\eta_0 + \theta \gamma)'T(x) - A(\eta_0 + \theta \gamma)}h(x)\\
    &= e^{\theta \cdot \gamma'T(x) - A(\eta_0 + \theta \gamma)}e^{\eta_0'T(x)}h(x)
\end{align*}
Thus, our new, lower dimensional natural parameter is $\theta$, our new sufficient statistic is $\gamma'T(x)$, and our new normalizing factor is $A(\eta_0 + \theta \gamma) = B(\theta)$. We also have $h(x) = e^{\eta_0'T(x)}h(x)$. \\

Also, interestingly, on $B$ which is a $1$-dimensional manifold, we can locally linearize so if we can restrict our family in question to a small subset, we see that it can be locally minimal in a sense. 

\subsection*{Completeness}

\begin{definition}[\textbf{Definition 7.5:} Completeness of a Sufficient Statistic]
    We say a sufficient statistic $T(X)$ is \emph{complete} for the family $\cal{P} = \{P_{\theta} \col \theta \in \Theta\}$ if $\E_{\theta}[f(T(X))] = 0$ for all $\theta$ implies that $f(T(X)) \overset{\tt{a.s.}}{=} 0$. 
\end{definition}

\begin{example}{Example}
    Consider the Laplace location family with minimal sufficient statistic $S(X) = (X_{(1)}, \ldots, X_{(n)})$ with $X_i \iidsim \tt{Lap}(\theta)$. Is $S(X)$ complete?\\

    The answer is no! Notice that $\bb{E}_{\theta}[X_{(n)} - X_{(1)}] = \bb{E}_{\theta}[Z_{(n)} - Z_{(1)}]$ where $Z_i = X_i - \theta \sim \tt{Lap}(0)$. Thus, we see that for 
    \[ f(S(X)) = (X_{(n)} - X_{(1)}) - \bb{E}[Z_{(n)} - Z_{(1)}] \]
    we have that
    \[ \bb{E}(f(S(X))) = 0 \]
    but $f(S(X)) \ne 0$ for all $\theta$!\\

    What this means is that $S(X)$, despite being sufficient, still has a lot of ``fluff''. It gives more information than needed in a sense. 
\end{example}